\section{Results and Discussion}\label{sec:mea-system-modeling-results}
% P:results-01
\Cref{tab:residual-summary} collects the deviations of the selected fit F1, the original Born fit F2 and the source-law fit F6 for every observation group; the Supplementary Material divides them by source and species.

\begin{table}[pos=!htb]
    \centering
    \small
    \caption{Cost and deviation summary for the two Born forms, each fitted to the 142-target objective unless stated otherwise.
    AARD is \(100\,n^{-1}\sum|y_{\mathrm{calc}}/y_{\mathrm{obs}}-1|\) over the stated group; the cost is the weighted least-squares cost \(C=\tfrac12\sum r_j^2\).
    Rows marked outside fit are tests outside the present (142-target, five-coordinate) fit, not data unseen during model development; canonical and transfer rows use input loadings corrected for a \(1\times10^{-4}\)~mol \COtwo/mol \MEA excess.
    The carbonate-share row covers nine Jakobsen points at 20 and \SI{40}{\degreeCelsius}, excluding the \SI{40}{\degreeCelsius} point at loading 0.21.}
    \label{tab:residual-summary}
    \begin{tabularx}{\linewidth}{@{}>{\raggedright\arraybackslash}X>{\raggedright\arraybackslash}p{3.1cm}r r r@{}}
        \toprule
        Group & Role & \(n\) & SSM+DS & Original Born \\
        \midrule
        \multicolumn{5}{@{}l}{\emph{Weighted cost}} \\
        Current objective & fitted & 142 & 32.991897 & 34.584032 \\
        Historical objective, historical fits & fitted & 160 & 47.625767 & 40.201537 \\
        Historical objective, current fits & re-scored & 160 & 53.769479 & 43.623963 \\
        \midrule
        \multicolumn{5}{@{}l}{\emph{AARD, \%}} \\
        Pressure, 40/\SI{60}{\degreeCelsius} & fitted & 48 & 13.53 & 17.31 \\
        Species, 20/40/\SI{60}{\degreeCelsius} & fitted & 94 & 8.28 & 8.23 \\
        Matin \ce{HCO3- + CO3^2-} pool & report-only & 18 & 59.81 & 35.12 \\
        Jou 1995 pressure, \SI{80}{\degreeCelsius} & outside fit & 11 & 9.24 & 11.21 \\
        Species, \SI{80}{\degreeCelsius} & outside fit & 11 & 11.15 & 7.36 \\
        Canonical pressure, \SI{80}{\degreeCelsius} & outside fit & 21 & 20.27 & 20.91 \\
        Canonical pressure, 40--\SI{80}{\degreeCelsius} & pooled & 104 & 20.58 & 20.78 \\
        Canonical pressure, 100--\SI{120}{\degreeCelsius} & extrapolation & 57 & 26.44 & 26.48 \\
        Pressure, 15 wt\% \MEA & transfer & 33 & 75.95 & 71.50 \\
        Pressure, 45 wt\% \MEA & transfer & 37 & 46.15 & 46.99 \\
        \midrule
        \multicolumn{5}{@{}l}{\emph{Calculated/observed ratio}} \\
        Jakobsen carbonate share, 20/\SI{40}{\degreeCelsius} & outside fit & 9 & 0.42--1.33 & 1.21--3.19 \\
        \bottomrule
    \end{tabularx}
\end{table}


\subsection{Calibrated model}\label{sec:results-calibration}
% P:results-02
The selected SSM+DS fit, F1, reproduces the 47 fitted \COtwo partial pressures at 40 and \SI{60}{\degreeCelsius} with an AARD of 13.98\,\% and the 94 fitted species at 20--\SI{60}{\degreeCelsius} with an AARD of 8.26\,\%, at a cost of 33.03.
The mean \(\ln(p_{\mathrm{calc}}/p_{\mathrm{obs}})\) of the fitted pressures is 0.012, so the calibration carries no material pooled pressure bias.
\Cref{fig:pressure,fig:speciation} compare the calculated pressures and species with all measurements, including those outside the fit.

\begin{figure}[pos=!htb]
    \centering
    \includegraphics[width=\linewidth]{pressure.pdf}
    \caption{F1 and the source-law fit F6 agree at the fitted 40 and \SI{60}{\degreeCelsius} isotherms, but F6 underestimates the measured \COtwo partial pressures above \SI{60}{\degreeCelsius}.
    \COtwo partial pressure against loading in 30 wt\% \MEA for 159 measured pressures from six studies, grouped by nominal temperature, and for 23 Wagner et al.\ pressures near \SI{353}{\kelvin} (c) and \SI{392}{\kelvin} (e).
    Open markers are measurements; dots (F1) and crosses (F6) are calculated at each measured state and are not interpolated.
    Only the 40 and \SI{60}{\degreeCelsius} panels contain fitted pressures.
    Three Xu and Rochelle rows above \SI{393.15}{\kelvin} are not evaluated and not shown.}
    \label{fig:pressure}
\end{figure}

\begin{figure}[pos=!htb]
    \centering
    \includegraphics[width=\linewidth]{speciation.pdf}
    \caption{F1 reproduces the fitted NMR and titration species at 20--\SI{60}{\degreeCelsius}, whereas the titration bicarbonate pool, which is not fitted, lies mostly above the calculated values.
    Liquid mole fraction against loading in 30 wt\% \MEA: (a) 76 titration values of Matin et al.\ at \SI{20}{\degreeCelsius}; (b)--(e) 55 NMR values of B\"ottinger et al., where \MEA and \MEAH are observed only as a sum.
    Circles are fitted measurements, squares are measurements outside the fit, and dots are F1 calculations at each measured state, joined to their measurement by a vertical segment.
    Bicarbonate values of both sources are compared with the modeled \ce{HCO3- + CO3^2-} pool.
    Matin's measurements, labeled \SI{21}{\degreeCelsius} in the source figures, are calculated at \SI{20}{\degreeCelsius}.}
    \label{fig:speciation}
\end{figure}

% P:results-03
At F1 the \HCOthree--water coefficient lies on its upper bound of \(+0.5\), with an outward cost gradient of \(-12.63\), so the cost would decrease further if the bound were relaxed; its value is set by the bound rather than estimated from the data.
For the four interior parameters, the conditional standard errors are 0.014 for \(k_{\MEACOO,\HtwoO}\), 0.017 for \(k_{\MEAH,\HtwoO}\) and 0.027 for \(k_{\MEAH,\MEACOO}\) (\cref{tab:binary-interaction-parameters}), and the largest correlation, 0.46, links \(k_{\MEAH,\HtwoO}\) with its slope \(b\).
The slope itself is \(b=\SI{-112.6}{\kelvin}\) with a conditional standard error of \SI{119.7}{\kelvin}, so these data, which constrain temperature only through pressures at 40 and \SI{60}{\degreeCelsius} and species at 20--\SI{60}{\degreeCelsius}, do not determine it.
It is retained so that all six fits share the same five parameters.
These standard errors rest on the declared residual scales and on residuals treated as independent (\cref{sec:data-methods}), and they are indicative only.

\subsection{Born term}\label{sec:results-born-off}
% P:results-04
Refitting all five parameters without the Born term (F3), within the same bounds and from two starts, gives a fitted-pressure AARD of 53.60\,\% and a fitted-species AARD of 10.79\,\%, against 13.98 and 8.26\,\% for F1.
The cost rises from 33.03 to 96.64, and the pressure targets carry 53.17 of this increase.
Both the \HCOthree--water and the \MEAH--\MEACOO coefficients reach their bounds, \(+0.5\) and \(-1\).
The tested refits within the declared bounds therefore do not recover the agreement of the model with a Born term; this result concerns these parameters and bounds and does not show that a Born contribution is physically necessary.

% P:results-05
A decomposition of the calculated pressure locates where the refit loses agreement.
Combining \COtwo phase equilibrium with the overall carbamate reaction R2 \(-\) R4 \(-\) R5, \(\ce{CO2 + 2 MEA <=> MEAH+ + MEACOO-}\), gives at every calculated state
\begin{equation}
    \ln p_{\COtwo}=S+Q-\ln K_{x,\mathrm{OV}}+H,
    \qquad
    S=\ln\frac{x_{\MEAH}\,x_{\MEACOO}}{x_{\MEA}^{2}},
    \label{eq:pressure-decomposition}
\end{equation}
where \(S\) is the speciation sum, \(Q=\sum_i\nu_{\mathrm{OV},i}\ln\gamma_i^{*}+\ln\gamma_{\COtwo}^{*}\) is the activity sum of the overall reaction, in which the \COtwo activity coefficient cancels, \(K_{x,\mathrm{OV}}\) is the mole-fraction equilibrium constant of the overall reaction, and \(H\) collects the vapor fugacity coefficient and reference fugacity of \COtwo.
At fixed speciation, liquid nonideality changes the \COtwo pressure only through \(Q\), and each Helmholtz contribution adds its own part to \(Q\).
For F3 relative to F1, \cref{eq:pressure-decomposition} splits the change of \(\ln p_{\COtwo}\) at each fitted pressure state into \(\Delta Q\), \(\Delta S\) and \(\Delta H\); the split reproduces the pressure-cost increase of 53.17 and closes within \(10^{-6}\) at every state (Supplementary Section~S3).

\begin{figure}[pos=!htb]
    \centering
    \includegraphics[width=\linewidth]{born-off-mechanism.pdf}
    \caption{Without the Born term, the refitted model F3 overestimates the fitted pressures at intermediate loadings, because the activity sum \(Q\) rises with loading faster than the speciation sum \(S\) falls.
    (a, b) \(\ln(p_{\mathrm{calc}}/p_{\mathrm{obs}})\) of the 47 fitted pressures for F1 and F3 in the nominal 40 and \SI{60}{\degreeCelsius} series; the shaded band is \(\pm0.3\), the declared pressure residual scale, not an acceptance interval.
    (c, d) Change from F1 to F3 of \(\ln p_{\COtwo}\) and of its contributions in \cref{eq:pressure-decomposition} at the same states.
    Points are calculated states at their measured temperatures and are not interpolated.}
    \label{fig:born-off}
\end{figure}

% P:results-06
Over the 47 fitted pressure states, the root-mean-square change of \(Q\), 1.09, exceeds that of \(S\), 0.82, so the refit changes the activity sum more than the speciation in aggregate (\cref{fig:born-off}).
\(\Delta Q\) rises with loading to about 2.3 at the highest loadings, where a nearly equal and opposite \(\Delta S\) cancels most of it; the net change of \(\ln p_{\COtwo}\) is largest at intermediate loadings, where F3 overestimates the measured pressures.
Two tests stated before the calculation locate the Born contribution of F1.
First, the Born part of the activity-coefficient sum \(\sum_i\nu_{\mathrm{OV},i}\ln\gamma_i^{*}\) of the overall reaction is negative, with a median of \(-0.91\) over the 32 fitted states at \SI{40}{\degreeCelsius}, and at the median-loading state its permittivity part, \(-0.29\), is less than half of the total, \(-0.88\); the hypothesis that the Born term raises this sum through the lower permittivity of \MEA--water is therefore rejected.
Second, the ion-concentration part of the Born contribution to \(Q\), defined as the Born contribution minus its value at the ion-free composition of the same state, falls steadily with loading: its Spearman correlation with loading is \(-0.997\) at \SI{40}{\degreeCelsius} and \(-0.996\) at \SI{60}{\degreeCelsius}, and its mean over the highest-loading third of the states is \(-1.76\) and \(-1.75\).
The non-Born activity changes of the refit offset the removed Born contribution fully or more in the lowest-loading third, with mean offset fractions of 1.00 and 1.58 at 40 and \SI{60}{\degreeCelsius}, but hardly at all in the highest-loading third, with 0.08 and \(-0.01\).
Within this model, the Born term therefore supplies a loading-dependent lowering of the activity sum that the five refitted interactions reproduce at low loading but not at high loading.
These are attributions within the calibrated model at its fitted states, under stated conventions for splitting the Born term; they are not measurements of single-ion activity coefficients.

\subsection{Born form and fitted observation set}\label{sec:results-born}
% P:results-07
The SSM+DS (F1) and original Born (F2) forms were fitted with the same parameters, bounds, weights, reaction constants and ion inputs, and they calibrate the data comparably.
With Matin's 18 pool targets not fitted, SSM+DS has the lower cost, 33.03 against 34.14, and the lower fitted-pressure AARD, 13.98 against 17.36\,\%, whereas original Born has the slightly lower fitted-species AARD, 8.16 against 8.26\,\% (\cref{tab:residual-summary}).
In F2 the \MEAH--\MEACOO coefficient also lies on its lower bound of \(-1\).
These are local cost comparisons, not statistical tests.

% P:results-08
Matin et al.\ obtained bicarbonate from a titration that assigns carbonate to bicarbonate, and they report bicarbonate above the NMR values of Jakobsen et al., whereas their free amine, protonated amine and carbamate agree more closely (Fig.~2, p.~6616; Fig.~3 and discussion, p.~6617)~\cite{matinFacileMethodDetermination2012,Jakobsen2005}.
On that basis the 18 Matin pool targets are not fitted in F1 and F2.
F1 lies below them, with an AARD of 59.59\,\% and a mean \(\ln(x_{\mathrm{calc}}/x_{\mathrm{obs}})\) of \(-1.06\), and F2 gives 34.99\,\%.

\begin{figure}[pos=!htb]
    \centering
    \includegraphics[width=\linewidth]{pool-effect.pdf}
    \caption{Fitting Matin's titration bicarbonate pool reverses the Born-form ranking.
    (a) \ce{HCO3- + CO3^2-} mole fraction at \SI{20}{\degreeCelsius}: Matin et al.\ measurements and the values calculated by F1 and F2, which do not fit the pool, and by F4 and F5, which fit it; the Matin state at loading 0.455 is outside both target sets.
    (b) Cost of the 141 base targets against the cost of the 18 pool targets; arrows join each Born form's fit without and with the pool.}
    \label{fig:pool-effect}
\end{figure}

% P:results-09
Adding the 18 pool targets to the objective reverses the order of the two forms (\cref{fig:pool-effect}): original Born then has the lower cost, 39.70 against 47.63 for SSM+DS.
Fitting the pool lowers its AARD to 42.57\,\% for SSM+DS (F4) and 19.56\,\% for original Born (F5), while the cost of the 141 base targets rises from 33.03 to 37.64 and from 34.14 to 35.91.
Original Born therefore moves toward the titration bicarbonate at a smaller cost to the other targets.
For these data, the pressure and speciation fits give no preference between the two Born formulations that is independent of the fitted observations, and independent solvation or activity data could help choose between them.
The residual signs do not establish which bicarbonate measurement is correct.

\subsection{Temperature and concentration outside the fit}\label{sec:results-80c}
% P:results-10
The R2 and R5 reaction shifts and the \COtwo dispersion energy of F1 and F2 were estimated with data that included \SI{80}{\degreeCelsius} (\cref{sec:data-methods}), so the comparisons at \SI{80}{\degreeCelsius} and above test conditions outside the five-parameter fit but are not independent temperature predictions.
On the 21 measured pressures at \SI{80}{\degreeCelsius}, F1 gives an AARD of 19.58\,\% with a mean \(\ln(p_{\mathrm{calc}}/p_{\mathrm{obs}})\) of \(-0.170\); the 11 Jou et al.\ measurements alone give 9.57\,\%.
For the 11 B\"ottinger species at \SI{80}{\degreeCelsius}, measured in a 31 wt\% solution, F1 gives 12.05\,\% and F2 8.08\,\%.
The 23 Wagner et al.\ pressures entered no fit or parameter selection; F1 gives 22.50\,\% with a mean \(\ln(p_{\mathrm{calc}}/p_{\mathrm{obs}})\) of \(-0.257\) on the 11 rows near \SI{353}{\kelvin} and 22.27\,\% with \(-0.265\) on the 12 rows near \SI{392}{\kelvin}, so it underestimates these pressures on average.
At 100--\SI{120}{\degreeCelsius}, the 54 evaluated pressures give 27.24\,\%.
These temperatures lie above the \SI{323.15}{\kelvin} upper limit of the R4 and R5 source correlations, so these values are extrapolations.
F2 gives similar values in each of these pressure groups (\cref{tab:residual-summary}).

% P:results-11
On the 105 measured pressures at 40--\SI{80}{\degreeCelsius}, F1 gives an AARD of 20.35\,\%.
This group combines fitted targets with rows outside the fit, so it measures pooled agreement, not accuracy outside the fit.
By source, the AARDs are 31.76\,\% for Aronu et al.\ (36 rows), 11.40\,\% for Hilliard (31 rows), 15.59\,\% for Jou et al.\ (28 rows) and 20.32\,\% for Idris et al.\ (10 rows)~\cite{Aronu2011,Hilliard2008,Jou1995,idrisSpeciationMEACO2Adducts2014}.
The mean \(\ln(p_{\mathrm{calc}}/p_{\mathrm{obs}})\) is \(-0.381\) for the Aronu rows and \(-0.136\) for all 105, so the pooled value is dominated by underestimation of the Aronu measurements.

% P:results-12
Applied unchanged at other amine concentrations, F1 gives AARDs of 76.71\,\% on 33 Aronu et al.\ pressures at 15 wt\% \MEA and 46.29\,\% on 37 at 45 wt\%, with mean \(\ln(p_{\mathrm{calc}}/p_{\mathrm{obs}})\) of \(+0.458\) and \(-0.522\), so the error changes sign with amine concentration and the parameters are calibrated for 30 wt\% \MEA.
Over nine Jakobsen NMR points at 20 and \SI{40}{\degreeCelsius}, the calculated carbonate share of dissolved carbon is 0.41--1.34 times the observed share for F1 and 1.19--3.21 times for F2~\cite{Jakobsen2005}.
The excluded point at \SI{40}{\degreeCelsius} and loading 0.21 has an observed share of 0.119, against 0.010--0.019 at the other four \SI{40}{\degreeCelsius} loadings, and the source dilution and loading definitions are unresolved; all ten comparisons are listed in the Supplementary Material.

\subsection{Source reaction laws}\label{sec:results-source-laws}
% P:results-13
F6 uses a different temperature description with the same 141 targets: R2 and R5 follow their source correlations, the \COtwo dispersion energy is the published \SI{169.21}{\kelvin}, and the five interaction parameters are refitted.
It fits the data at or below \SI{60}{\degreeCelsius} as well as F1, with a cost of 33.06 against 33.03, but represents higher temperatures worse (\cref{fig:pressure}).
On the 21 measured pressures at \SI{80}{\degreeCelsius} it gives 28.91\,\% against 19.58\,\%, on the Wagner rows near \SI{353}{\kelvin} 51.60\,\% against 22.50\,\%, and on the 54 evaluated pressures at 100--\SI{120}{\degreeCelsius} 51.05\,\% against 27.24\,\%, with a mean \(\ln(p_{\mathrm{calc}}/p_{\mathrm{obs}})\) of \(-0.731\) on the last group, so its calculated pressures are too low.
At 15 wt\% \MEA, F6 gives 31.12\,\% against 76.71\,\% for F1; at 45 wt\% it gives 55.18\,\% against 46.29\,\%.

% P:results-14
F6 changes the reaction laws, the \COtwo dispersion energy and the refitted interactions together, so the difference from F1 does not isolate the effect of the reaction laws.
Both fits use their reaction correlations over the same declared domain, 293.15--\SI{393.15}{\kelvin}, which extends beyond the source ranges; the R4 source correlation, for example, ends at \SI{323.15}{\kelvin}.
The comparison therefore assesses extrapolated model behavior.
Data at or below \SI{60}{\degreeCelsius} do not distinguish between the two temperature descriptions, and the higher-temperature agreement of F1 rests on shifts estimated with \SI{80}{\degreeCelsius} data.
For modelers, the temperature response of this model must come from data that carry it directly, such as heat of absorption or equilibrium measurements above \SI{60}{\degreeCelsius}.

\subsection{Comparison with published models}\label{sec:results-literature}
% P:results-15
Like the electrolyte NRTL model of Zhang et al., which was fitted to \COtwo solubility, heat of absorption, heat capacity and NMR speciation~\cite{Zhang2011}, and unlike that of Akula et al., which left NMR speciation outside its fit~\cite{Akula2023a}, the present model fits speciation together with pressure, so the residuals of its 94 fitted species targets are calibration results (\cref{tab:literature-model-comparison}).
Unlike the ePC-SAFT models of MDEA~\cite{Uyan2015,Wangler2018,Bulow2021a}, it fits ion interactions to the loaded ternary mixture, including the \MEAH--\MEACOO interaction, which exists only in loaded solution.
Published pressure deviations for 30 wt\% \MEA cover different row sets and temperature ranges, so no accuracy ranking against those models is claimed.

\subsection{Scope and limitations}\label{sec:results-limits}
% P:results-16
F1 describes \COtwo partial pressure in 30 wt\% \MEA, fitted at 40 and \SI{60}{\degreeCelsius}, and liquid speciation, fitted at 20--\SI{60}{\degreeCelsius}; its comparisons at 80--\SI{120}{\degreeCelsius} rest partly on reaction shifts estimated with \SI{80}{\degreeCelsius} data.
Its ion parameters are effective values: the \MEAH and \MEACOO parameters were fitted to speciation data, the \HCOthree--water coefficient lies on its bound, and the \MEAH--water slope is not determined.
The conditional standard errors treat residuals as independent and exclude the uncertainty of the fixed inputs.
The Matin values are a transcription not checked against the source's supporting tables.

% P:results-17
Liquid density is not fitted.
At an assumed pressure of \SI{101325}{\pascal}, which the source does not report, F1 overpredicts the four loaded densities of Amundsen et al.\ at 50 and \SI{70}{\degreeCelsius} and loadings 0.3 and 0.4~mol \COtwo/mol \MEA, on a \COtwo-free 30 wt\% \MEA solvent basis, by 13.0--17.8\,\%~\cite[Table~3, p.~3097]{Amundsen2009}.
The source gives a general uncertainty estimate of \(\pm\SI{0.002}{\gram\per\cubic\centi\metre}\) for loaded solutions but excludes the highest temperatures and loadings from it~\cite[pp.~3099--3100]{Amundsen2009}, so these deviations are not weighted by uncertainty.
Heat capacity and heat of absorption were not evaluated for F1.
The parameter set is therefore a phase- and speciation-equilibrium model; absorber and stripper calculations would require density and caloric qualification first.
